\documentclass[aspectratio=169,10pt]{beamer}

% --- Paquetes Básicos ---
\usepackage[utf8]{inputenc}
\usepackage[spanish]{babel}
\usepackage{graphicx}
\usepackage{booktabs}
\usepackage{amsmath}
\usepackage{amsfonts}
\usepackage{tikz}
\usetikzlibrary{shapes.geometric, arrows.meta, positioning, fit, backgrounds}

% --- Tema y Colores Premium UTEC Posgrado ---
\usetheme{Madrid}
\usecolortheme{whale}

\definecolor{UtecCyan}{RGB}{0, 191, 255}       % Azul cian brillante
\definecolor{UtecDarkGray}{RGB}{50, 50, 50}    % Gris Oscuro
\definecolor{LightGray}{RGB}{245, 245, 245}

\setbeamercolor{palette primary}{bg=UtecDarkGray, fg=white}
\setbeamercolor{palette secondary}{bg=UtecDarkGray!85, fg=white}
\setbeamercolor{palette tertiary}{bg=UtecCyan, fg=UtecDarkGray}
\setbeamercolor{titlelike}{parent=palette primary, fg=UtecDarkGray}
\setbeamercolor{structure}{fg=UtecCyan}
\setbeamercolor{normal text}{fg=black}

\setbeamertemplate{navigation symbols}{}
\setbeamertemplate{footline}[frame number]

% --- Configuración de Estilos TikZ para Beamer ---
\tikzset{
    block/.style={draw, rectangle, rounded corners, fill=LightGray, text width=2.6cm, align=center, minimum height=0.6cm, scale=0.6, thick},
    cloud/.style={draw, ellipse, fill=UtecCyan!20, text width=2.4cm, align=center, minimum height=0.6cm, scale=0.6, thick},
    decision/.style={draw, diamond, aspect=2, fill=orange!15, text width=2.4cm, align=center, minimum height=0.6cm, scale=0.6, thick},
    line/.style={draw, -{Latex[scale=0.8]}, thick, scale=0.6}
}

% --- Metadatos ---
\title[Sustentación de Arquitectura de Datos]{INGENIERÍA DE DATOS -- PROYECTO FINAL}
\subtitle{Defensa de Arquitectura Medallion Resiliente, Data Mesh y Gobernanza}
\author[Grupo G2]{Integrantes: \\ \textbf{Grupo G2 (Sección 1 - g102)}}
\institute[UTEC]{Universidad de Ingeniería y Tecnología - UTEC Posgrado}
\date{18 de Julio de 2026}

\begin{document}

% --- Slide 1: Portada ---
{
\setbeamercolor{background canvas}{bg=UtecDarkGray}
\begin{frame}[plain]
    \centering
    \vspace{1.2cm}
    {\large \color{UtecCyan} \textbf{UTEC POSGRADO -- MAESTRÍA EN INGENIERÍA DE DATOS}} \\[0.2cm]
    {\LARGE \color{white} \textbf{Sustentación de Arquitectura y Gobernanza Técnica}} \\[0.4cm]
    {\Large \color{UtecCyan} \textbf{Pipeline Medallion Resiliente y Descentralizado}} \\[0.8cm]
    \color{white}
    \footnotesize
    Compañía Minera Los Andes S.A. | Control de Inventario por Consignación \\
    Grupo G2 (g102) -- Sustentación de Examen Final \\
    Julio 2026
\end{frame}
}

% --- Slide 2: Arquitectura del Pipeline (TikZ) ---
\begin{frame}{Arquitectura del Pipeline Medallion}
    \centering
    \begin{tikzpicture}[scale=0.85, every node/.style={scale=0.85}, node distance=1.2cm, auto, >=latex', thick]
        % Capas Medallion
        \node[draw, cylinder, fill=LightGray, text=UtecDarkGray, align=center, scale=0.7] (raw) {ADLS Gen2 Raw\\(CSV SAP / Maestro)};
        \node[draw, rectangle, rounded corners, fill=UtecDarkGray, text=white, right=1.2cm of raw, align=center, text width=2.4cm, scale=0.7] (bronze) {\textbf{Capa Bronze}\\Inmutable Delta};
        \node[draw, rectangle, rounded corners, fill=UtecDarkGray, text=white, right=1.2cm of bronze, align=center, text width=2.4cm, scale=0.7] (silver) {\textbf{Capa Silver}\\Limpieza/Enmascarado};
        \node[draw, rectangle, rounded corners, fill=UtecDarkGray, text=white, right=1.2cm of silver, align=center, text width=2.4cm, scale=0.7] (gold) {\textbf{Capa Gold}\\3 Data Products};
        
        \node[draw, rectangle, fill=red!15, text=red, below=0.8cm of silver, align=center, scale=0.7] (quarantine) {Cuarentena Delta\\(Registros Errantes)};
        
        % Flechas de Flujo
        \draw[->, double, draw=UtecDarkGray] (raw) -- node[above, scale=0.6]{abfss / SSL} (bronze);
        \draw[->, draw=UtecDarkGray] (bronze) -- node[above, scale=0.6]{try\_cast / join} (silver);
        \draw[->, draw=red, dashed] (silver) -- node[left, scale=0.6]{Rechazos} (quarantine);
        \draw[->, draw=UtecDarkGray] (silver) -- node[above, scale=0.6]{SQL Views} (gold);
    \end{tikzpicture}
    \vspace{0.3cm}
    \footnotesize
    Todo el pipeline se ejecuta bajo el metastore de \textbf{Unity Catalog}, garantizando el linaje de datos de extremo a extremo y auditoría integrada.
\end{frame}

% --- Slide 3: El Ecosistema Data Mesh ---
\begin{frame}{El Ecosistema de Data Products (Data Mesh)}
    \begin{columns}
        \begin{column}{0.5\textwidth}
            \textbf{Interdependencia Semántica de Productos:}
            \begin{itemize}
                \setlength{\itemsep}{0.4em}
                \item \textbf{Capa Silver Unificada:} Los 3 productos en Gold leen de la tabla común \texttt{silver\_movimientos\_consignacion}.
                \item \textbf{Retroalimentación Operativa (DP2 $\rightarrow$ DP1):} El análisis de consumo mensual real (DP2) permite calibrar y reajustar los stocks de seguridad y puntos de reorden en el maestro de materiales, recalculando dinámicamente las alertas de compra (DP1).
                \item \textbf{Auditoría de Pagos (DP2 $\leftrightarrow$ DP3):} Cada USD reportado como costo de consumo en DP2 debe corresponder a una orden de pago pre-aprobada y conciliada en DP3.
            \end{itemize}
        \end{column}
        \begin{column}{0.48\textwidth}
            \centering
            \begin{tikzpicture}[scale=0.75, every node/.style={scale=0.75}, node distance=1.1cm, auto, >=latex', thick]
                \node[draw, rectangle, fill=blue!15, text width=5cm, align=center, scale=0.75] (silver) {\textbf{Capa Silver Unificada}\\ \texttt{silver\_movimientos\_consignacion}};
                \node[draw, rectangle, fill=orange!15, below left=1.2cm and 0.5cm of silver, text width=2.8cm, align=center, scale=0.75] (dp1) {\textbf{DP1: Alertas}\\ \texttt{mv\_alerts}};
                \node[draw, rectangle, fill=orange!15, below=1.2cm of silver, text width=2.8cm, align=center, scale=0.75] (dp2) {\textbf{DP2: Consumos}\\ \texttt{mv\_consumo}};
                \node[draw, rectangle, fill=orange!15, below right=1.2cm and 0.5cm of silver, text width=2.8cm, align=center, scale=0.75] (dp3) {\textbf{DP3: Pagos}\\ \texttt{mv\_conciliation}};
                
                \draw[->] (silver) -- (dp1);
                \draw[->] (silver) -- (dp2);
                \draw[->] (silver) -- (dp3);
                \draw[->, dashed, draw=red, <->] (dp2) -- (dp3);
                \draw[->, dashed, draw=blue] (dp2) [bend left=45] to (dp1);
            \end{tikzpicture}
        \end{column}
    \end{columns}
\end{frame}

% --- Slide 4: Capa Bronze ---
\begin{frame}{Capa Bronze: Ingesta Inmutable}
    \begin{columns}
        \begin{column}{0.5\textwidth}
            \textbf{Notebook 1 - Flujo de Ingesta:}
            \begin{itemize}
                \setlength{\itemsep}{0.4em}
                \item \textbf{Conectividad ADLS Gen2:} Lectura directa mediante protocolo de red seguro \texttt{abfss} usando llaves SAS.
                \item \textbf{Mitigación del Delimitador:} Configuración del lector con delimitador ``;`` de SAP para evitar que Spark lea columnas agrupadas.
                \item \textbf{Estandarización Delta:} Reemplazo sistemático de espacios y puntos de SAP en nombres de columna por guiones bajos.
            \end{itemize}
        \end{column}
        \begin{column}{0.48\textwidth}
            \centering
            \begin{tikzpicture}[node distance=0.8cm, auto, >=latex', thick]
                \node[cloud] (start) {Inicio};
                \node[block, below=of start] (step1) {Cargar Llaves SAS};
                \node[block, below=of step1] (step2) {Leer CSVs SAP (Delimiter ;)};
                \node[block, below=of step2] (step3) {Añadir Timestamp e Ingested At};
                \node[block, below=of step3] (step4) {Normalizar Nombres de Columnas};
                \node[block, below=of step4] (step5) {Guardar Delta Bronze};
                \node[cloud, below=of step5] (end) {Fin};
                
                \draw[line] (start) -- (step1);
                \draw[line] (step1) -- (step2);
                \draw[line] (step2) -- (step3);
                \draw[line] (step3) -- (step4);
                \draw[line] (step4) -- (step5);
                \draw[line] (step5) -- (end);
            \end{tikzpicture}
        \end{column}
    \end{columns}
\end{frame}

% --- Slide 5: Capa Silver ---
\begin{frame}{Capa Silver: Calidad y Enriquecimiento}
    \begin{columns}
        \begin{column}{0.5\textwidth}
            \textbf{Notebook 2 - Procesamiento y Cuarentena:}
            \begin{itemize}
                \setlength{\itemsep}{0.3em}
                \item \textbf{Tolerancia a Fechas:} Intento de casteo a múltiples formatos de SAP (ISO y local regional) mediante \texttt{coalesce}.
                \item \textbf{Casteo Resiliente:} Uso de \texttt{try\_cast} para evitar excepciones por registros desalineados.
                \item \textbf{Enriquecimiento:} Cruce con maestro de materiales y enmascaramiento PII del usuario de almacén.
                \item \textbf{Cuarentena:} Desvío automático de errores lógicos y cantidades nulas/cero.
            \end{itemize}
        \end{column}
        \begin{column}{0.48\textwidth}
            \centering
            \begin{tikzpicture}[node distance=0.6cm, auto, >=latex', thick]
                \node[cloud] (start) {Inicio};
                \node[block, below=of start] (read) {Leer tablas Bronze};
                \node[block, below=of read] (transform) {Aplicar try\_to\_date y try\_cast};
                \node[block, below=of transform] (join) {Cruce con Maestro};
                \node[decision, below=of join] (quality) {¿Calidad OK?};
                \node[block, below left=0.6cm and 0.1cm of quality, fill=green!10, text width=2.2cm] (clean) {Enmascarar PII y Limpiar};
                \node[block, below right=0.6cm and 0.1cm of quality, fill=red!10, text width=2.2cm] (quarantine) {Desviar a Cuarentena};
                \node[block, below=1.8cm of quality] (write) {Escribir Tablas Silver};
                \node[cloud, below=of write] (end) {Fin};
                
                \draw[line] (start) -- (read);
                \draw[line] (read) -- (transform);
                \draw[line] (transform) -- (join);
                \draw[line] (join) -- (quality);
                \draw[line] (quality) -| (clean);
                \draw[line] (quality) -| (quarantine);
                \draw[line] (clean) |- (write);
                \draw[line] (quarantine) |- (write);
                \draw[line] (write) -- (end);
            \end{tikzpicture}
        \end{column}
    \end{columns}
\end{frame}

% --- Slide 6: Capa Gold ---
\begin{frame}{Capa Gold: Creación de Data Products}
    \begin{columns}
        \begin{column}{0.5\textwidth}
            \textbf{Notebook 3 - Modelamiento Semántico:}
            \begin{itemize}
                \setlength{\itemsep}{0.4em}
                \item \textbf{Cómputo en el Cluster:} Reemplazo de vistas materializadas por vistas estándar debido a restricciones de Warehouse SQL Pro/Serverless en el clúster interactivo.
                \item \textbf{DP1 (Alertas):} Alertas dinámicas al caer el stock por debajo del stock de seguridad o reorden.
                \item \textbf{DP2 (Consumos):} Consumos netos en USD mensuales por categoría de repuesto.
                \item \textbf{DP3 (Conciliación):} Consolidación contable mensual agrupada por Documento de Material SAP.
            \end{itemize}
        \end{column}
        \begin{column}{0.48\textwidth}
            \centering
            \begin{tikzpicture}[node distance=0.8cm, auto, >=latex', thick]
                \node[cloud] (start) {Inicio};
                \node[block, below=of start] (step1) {Verificar Esquemas Gold};
                \node[block, below=of step1] (step2) {Calcular Alertas (DP1)};
                \node[block, below=of step2] (step3) {Calcular Consumos (DP2)};
                \node[block, below=of step3] (step4) {Calcular Pagos Conciliados (DP3)};
                \node[block, below=of step4] (step5) {Crear SQL Views en Gold};
                \node[cloud, below=of step5] (end) {Fin};
                
                \draw[line] (start) -- (step1);
                \draw[line] (step1) -- (step2);
                \draw[line] (step2) -- (step3);
                \draw[line] (step3) -- (step4);
                \draw[line] (step4) -- (step5);
                \draw[line] (step5) -- (end);
            \end{tikzpicture}
        \end{column}
    \end{columns}
\end{frame}

% --- Slide 7: Capa Gold Gobernanza ---
\begin{frame}{Gobernanza y Data Contracts}
    \begin{columns}
        \begin{column}{0.5\textwidth}
            \textbf{Notebook 4 - Gobernanza y Trazabilidad:}
            \begin{itemize}
                \setlength{\itemsep}{0.4em}
                \item \textbf{Diccionario de Datos:} Inyección de metadatos de columnas directamente en el metastore con comandos SQL \texttt{COMMENT ON}.
                \item \textbf{Data Contracts (YAML):} Función automatizada en Python que lee el esquema Delta de Gold y genera un contrato en YAML para los consumidores descentralizados.
                \item \textbf{Monitoreo de SLAs:} Vista de auditoría mockeada debido a la falta de permisos de cuenta para consultar \texttt{system.query.history}.
            \end{itemize}
        \end{column}
        \begin{column}{0.48\textwidth}
            \centering
            \begin{tikzpicture}[node distance=0.8cm, auto, >=latex', thick]
                \node[cloud] (start) {Inicio};
                \node[block, below=of start] (step1) {Registrar COMMENTs en Columnas};
                \node[block, below=of step1] (step2) {Crear Vista SLAs (Audit)};
                \node[block, below=of step2] (step3) {Extraer Esquema de Gold en Python};
                \node[block, below=of step3] (step4) {Generar Diccionario YAML};
                \node[block, below=of step4] (step5) {Publicar Contrato de Datos};
                \node[cloud, below=of step5] (end) {Fin};
                
                \draw[line] (start) -- (step1);
                \draw[line] (step1) -- (step2);
                \draw[line] (step2) -- (step3);
                \draw[line] (step3) -- (step4);
                \draw[line] (step4) -- (step5);
                \draw[line] (step5) -- (end);
            \end{tikzpicture}
        \end{column}
    \end{columns}
\end{frame}

% --- Slide 8: Desafíos Técnicos de SAP y Soluciones (Slide Unificado) ---
\begin{frame}{Desafíos Técnicos de SAP y Soluciones de Ingeniería}
    \small
    \begin{table}[h]
        \begin{tabular}{lll}
            \toprule
            \textbf{Dolor Técnico / Síntoma} & \textbf{Causa Raíz Operativa} & \textbf{Solución de Arquitectura Implementada} \\
            \midrule
            \rowcolor{red!10} \textbf{233.69K en Cuarentena} & Ceros a la izquierda en IDs SAP (`000...10`) & Estandarizar con \texttt{ltrim(col, '0')} antes de unir \\
            & vs maestro sin ceros (`10`). Fallo del Join. & en la capa Silver. \\
            \midrule
            \rowcolor{red!10} \textbf{Stocks Acumulados Negativos} & Falta de fotos de saldo de apertura del ERP. & Inyectar saldos iniciales (Mov. 561 de SAP) \\
            & Los consumos iniciales restaban de cero. & en la víspera en Bronze/Silver. \\
            \midrule
            \rowcolor{red!10} \textbf{Casteo Numérico Fallido} & Formato decimal de SAP con comas y puntos. & Reemplazar comas por puntos en Silver con \\
            & PySpark retorna NULL en try\_cast normal. & regex (\texttt{regexp\_replace}) antes del cast. \\
            \midrule
            \rowcolor{red!10} \textbf{Restricciones de Permisos} & Clúster compartido bloquea vistas materializadas & Usar vistas estándar en Gold y mockear reporte \\
            \textbf{en la Nube} & y lectura del catálogo \texttt{system.query.history}. & de SLAs en base a consultas simuladas. \\
            \bottomrule
        \end{tabular}
    \end{table}
\end{frame}

% --- Slide 9: Seguridad PII y SLAs ---
\begin{frame}{Seguridad PII y SLAs}
    \begin{columns}
        \begin{column}{0.5\textwidth}
            \textbf{Enmascaramiento Dinámico de Operadores (PII):}
            \begin{itemize}
                \setlength{\itemsep}{0.4em}
                \item \textbf{Problema:} Nombre de usuario de almacén de SAP es dato sensible.
                \item \textbf{Solución:} Creamos la UDF SQL \texttt{pii\_mask\_user} en Unity Catalog.
                \item \textbf{Gobernanza:} Uso de \texttt{is\_member('utec-admin')} para ocultar la información en tiempo de ejecución a analistas comunes.
            \end{itemize}
        \end{column}
        \begin{column}{0.48\textwidth}
            \textbf{Métricas de SLA de los Data Products:}
            \begin{itemize}
                \setlength{\itemsep}{0.4em}
                \item \textbf{Latencia Promedio:} Consultas interactivas completadas en $<2.5$ segundos (SLA cumplido).
                \item \textbf{Porcentaje de Éxito:} Mayor al $99\%$ de consultas exitosas.
                \item \textbf{Disponibilidad:} Datos actualizados diariamente listos antes del inicio de la jornada en mina.
            \end{itemize}
        \end{column}
    \end{columns}
\end{frame}

% --- Slide 10: Orquestación (Workflows) ---
\begin{frame}{Orquestación de Pipelines en la Nube}
    \begin{columns}
        \begin{column}{0.5\textwidth}
            \textbf{El Job de Producción (Databricks Workflows):}
            \begin{itemize}
                \setlength{\itemsep}{0.4em}
                \item \textbf{Secuenciación de Tareas:} Ingesta Bronze $\rightarrow$ Calidad y Enriquecimiento Silver $\rightarrow$ Vistas Gold $\rightarrow$ Gobernanza y Data Contracts.
                \item \textbf{Programación:} Ejecución diaria programada a las **01:00 AM** (Lima).
                \item \textbf{Alertas de Fallo:} Alertas por correo electrónico integradas nativamente en caso de caída de tareas.
            \end{itemize}
        \end{column}
        \begin{column}{0.45\textwidth}
            \centering
            \begin{beamercolorbox}[shadow=true,rounded=true,wd=\textwidth]{palette secondary}
                \centering \textbf{Flujo del Workflow}
            \end{beamercolorbox}
            \vspace{0.4cm}
            \small
            \begin{tikzpicture}[scale=0.75, every node/.style={scale=0.75}, node distance=1.1cm, >=latex', thick]
                \node[draw, rectangle, fill=LightGray, align=center, text width=4.5cm] (t1) {Task 1: Ingesta Bronze};
                \node[draw, rectangle, fill=LightGray, align=center, text width=4.5cm, below=0.6cm of t1] (t2) {Task 2: Procesamiento Silver};
                \node[draw, rectangle, fill=LightGray, align=center, text width=4.5cm, below=0.6cm of t2] (t3) {Task 3: Vistas Gold};
                \node[draw, rectangle, fill=LightGray, align=center, text width=4.5cm, below=0.6cm of t3] (t4) {Task 4: Gobernanza / YAML};
                
                \draw[->, draw=UtecDarkGray] (t1) -- (t2);
                \draw[->, draw=UtecDarkGray] (t2) -- (t3);
                \draw[->, draw=UtecDarkGray] (t3) -- (t4);
            \end{tikzpicture}
        \end{column}
    \end{columns}
\end{frame}

% --- Slide 11: Oportunidades de Mejora de Arquitectura ---
\begin{frame}{Oportunidades de Mejora de Arquitectura}
    \begin{itemize}
        \setlength{\itemsep}{0.6em}
        \item \textbf{Databricks Auto Loader:} Reemplazar la ingesta programada por un flujo continuo Near Real-Time en Bronze usando notificaciones de eventos (Azure Event Grid).
        \item \textbf{Orquestación Híbrida Desacoplada:} Mantener Apache Airflow en local como la torre de control de negocio, gatillando de forma remota el Job en Databricks mediante la API, evitando problemas de memoria locales de la VM.
        \item \textbf{Data Mesh Real (Multi-Workspace):} Aislar los recursos de Logística y Finanzas en workspaces separados de Databricks para descentralizar costos y gobernar mediante compartición segura con **Delta Sharing**.
        \item \textbf{Trazabilidad Unificada:} Implementar **OpenLineage** de principio a fin, conectando la traza desde el archivo crudo, Airflow, Databricks, hasta el dashboard de BI.
    \end{itemize}
\end{frame}

% --- Slide 12: Fin ---
{
\setbeamercolor{background canvas}{bg=UtecDarkGray}
\begin{frame}[plain]
    \centering
    \vspace{1.5cm}
    {\Huge \color{UtecCyan} \textbf{Gracias}} \\[0.8cm]
    {\large \color{white} \textbf{¿Preguntas sobre la arquitectura o gobernanza técnica?}} \\[1cm]
    \footnotesize
    \color{white!70}
    Grupo G2 (Sección 1 - g102) | Maestría en Ingeniería de Datos - UTEC
\end{frame}
}

\end{document}
